% ============================================================
\chapter{Quasi-m\'etriques et Topologie Asym\'etrique}
\label{ch:quasi_metriques}
% ============================================================

Depuis Fr\'echet (1906), la distance v\'erifie trois axiomes~: positivit\'e, sym\'etrie, in\'egalit\'e triangulaire. Mais de nombreuses situations r\'eelles sont intrins\`equement asym\'etriques~: le temps de trajet en mont\'ee diff\`ere de celui en descente, le co\^ut de calcul d'une approximation d\'epend du sens du raffinement, et en informatique th\'eorique, la distance de Smyth entre programmes mesure un rapprochement qui n'est pas r\'eciproque. Abandonner l'axiome de sym\'etrie conduit aux \emph{quasi-m\'etriques}, un cadre d\'evelopp\'e notamment par Hans-Peter K\"unzi \`a partir des ann\'ees 1980, qui r\'ev\`ele une topologie riche et surprenante~: chaque quasi-m\'etrique engendre \emph{deux} topologies, une vers l'avant et une vers l'arri\`ere.

\begin{intuition}
Les quasi-m\'etriques g\'en\'eralisent les m\'etriques en abandonnant l'axiome de
sym\'etrie~: la <<~distance~>> de $x$ \`a $y$ peut diff\'erer de celle de $y$ \`a $x$.
Cette asym\'etrie, loin d'\^etre une pathologie, capture naturellement des notions de co\^ut
orient\'e, de raffinement, et de complexit\'e en informatique.
\end{intuition}

% ------------------------------------------------------------
\section{D\'efinitions fondamentales}
% ------------------------------------------------------------

\begin{definition}[Quasi-m\'etrique]
Une \emph{quasi-m\'etrique} sur un ensemble $X$ est une fonction
$d\colon X \times X \to [0,+\infty]$ satisfaisant~:
\begin{enumerate}
  \item $d(x,x) = 0$ pour tout $x \in X$;
  \item $d(x,z) \leq d(x,y) + d(y,z)$ pour tous $x,y,z \in X$ (in\'egalit\'e triangulaire).
\end{enumerate}
Si de plus $d(x,y) = 0 = d(y,x) \Rightarrow x = y$, on dit que $d$ est \emph{$T_0$-s\'epar\'ee}.
Le couple $(X,d)$ est un \emph{espace quasi-m\'etrique}.
\end{definition}

\begin{remarque}
Nous autorisons $d$ \`a prendre la valeur $+\infty$, ce qui est essentiel pour certaines
constructions (produits infinis, domaines de complexit\'e).
\end{remarque}

\begin{definition}[Quasi-m\'etrique conjugu\'ee et sym\'etris\'ee]
Soit $(X,d)$ un espace quasi-m\'etrique.
\begin{itemize}
  \item La \emph{conjugu\'ee} de $d$ est $d^{-1}(x,y) = d(y,x)$.
  \item La \emph{sym\'etris\'ee} de $d$ est $d^s(x,y) = \max(d(x,y), d^{-1}(x,y))$.
  \item La \emph{somme sym\'etris\'ee} est $\hat{d}(x,y) = d(x,y) + d^{-1}(x,y)$.
\end{itemize}
La sym\'etris\'ee $d^s$ est une m\'etrique (\'eventuellement \`a valeurs dans $[0,+\infty]$)
d\`es que $d$ est $T_0$-s\'epar\'ee.
\end{definition}

% ------------------------------------------------------------
\section{Topologie associ\'ee \`a une quasi-m\'etrique}
% ------------------------------------------------------------

\begin{definition}[Boules ouvertes]
Pour $(X,d)$ quasi-m\'etrique, $x \in X$ et $\varepsilon > 0$~:
\[
  B_d(x,\varepsilon) = \{y \in X : d(x,y) < \varepsilon\}
\]
est la \emph{boule ouverte avant} (ou \emph{boule directe}).
La \emph{topologie directe} $\tau_d$ est la topologie engendr\'ee par ces boules.
\end{definition}

\begin{proposition}
Les boules ouvertes $B_d(x,\varepsilon)$ forment une base de la topologie $\tau_d$.
L'espace $(X, \tau_d)$ est $T_0$ si et seulement si $d$ est $T_0$-s\'epar\'ee.
\end{proposition}

\begin{proof}
Soient $y \in B_d(x_1,\varepsilon_1) \cap B_d(x_2,\varepsilon_2)$. Posons
$\delta = \min(\varepsilon_1 - d(x_1,y), \varepsilon_2 - d(x_2,y)) > 0$. Pour tout
$z \in B_d(y,\delta)$, on a $d(x_i,z) \leq d(x_i,y) + d(y,z) < d(x_i,y) + \delta \leq
\varepsilon_i$. Donc $B_d(y,\delta) \subseteq B_d(x_1,\varepsilon_1) \cap
B_d(x_2,\varepsilon_2)$, ce qui montre que les boules forment une base.

Pour la s\'eparation $T_0$~: si $\tau_d$ est $T_0$, alors pour $x \neq y$, il existe
$\varepsilon > 0$ tel que, disons, $y \notin B_d(x,\varepsilon)$, d'o\`u $d(x,y) \geq
\varepsilon > 0$. Comme l'un des deux ($d(x,y)$ ou $d(y,x)$) est non nul, $d$ est
$T_0$-s\'epar\'ee.
\end{proof}

\begin{definition}[Ordre de sp\'ecialisation quasi-m\'etrique]
L'ordre de sp\'ecialisation de $\tau_d$ est~:
\[
  x \leq_d y \quad \Longleftrightarrow \quad d(x,y) = 0.
\]
\end{definition}

\begin{proposition}
L'ordre $\leq_d$ est un pr\'eordre (ordre partiel si $d$ est $T_0$-s\'epar\'ee), et
pour $d$ $T_0$-s\'epar\'ee, $d(x,y) = 0$ \'equivaut \`a <<~$y$ est dans tout ouvert
contenant $x$~>>.
\end{proposition}

% ------------------------------------------------------------
\section{Exemples fondamentaux}
% ------------------------------------------------------------

\begin{exemple}[Quasi-m\'etrique de Sorgenfrey]
Sur $\R$, posons~:
\[
  d_S(x,y) = \begin{cases} y - x & \text{si } y \geq x, \\ 1 & \text{si } y < x. \end{cases}
\]
Alors $\tau_{d_S}$ est la topologie de Sorgenfrey (base de semi-intervalles $[a,b)$).
\end{exemple}

\begin{exemple}[Quasi-m\'etrique de Baire]
Sur $\N^\N$ (suites infinies de naturels), d\'efinissons~:
\[
  d(f,g) = \begin{cases} 0 & \text{si } f = g, \\
  2^{-\min\{n : f(n) \neq g(n)\}} & \text{si } f \neq g \text{ et } f \sqsubseteq g, \\
  1 & \text{sinon,} \end{cases}
\]
o\`u $f \sqsubseteq g$ signifie que $f$ est un pr\'efixe de $g$ (au sens o\`u $f(k) \leq g(k)$
pour les premiers indices). Cette quasi-m\'etrique capture l'id\'ee d'approximation~: $d(f,g)$
est petit si $g$ \'etend $f$ sur un long pr\'efixe.
\end{exemple}

\begin{exemple}[Quasi-m\'etrique de complexit\'e]
Soit $\Sigma^*$ l'ensemble des mots sur un alphabet $\Sigma$. D\'efinissons~:
\[
  d(u,v) = \begin{cases} 0 & \text{si } u \text{ est pr\'efixe de } v, \\
  2^{-|u \wedge v|} & \text{sinon,} \end{cases}
\]
o\`u $|u \wedge v|$ est la longueur du plus long pr\'efixe commun. C'est une
quasi-m\'etrique $T_0$-s\'epar\'ee dont la topologie directe est la topologie de Scott sur
les mots (partiels).
\end{exemple}

% ------------------------------------------------------------
\section{Quasi-uniformit\'es}
% ------------------------------------------------------------

\begin{definition}[Quasi-uniformit\'e]
Une \emph{quasi-uniformit\'e} sur $X$ est un filtre $\mathcal{U}$ de sous-ensembles de
$X \times X$ tel que~:
\begin{enumerate}
  \item Pour tout $U \in \mathcal{U}$, la diagonale $\Delta_X \subseteq U$.
  \item Pour tout $U \in \mathcal{U}$, il existe $V \in \mathcal{U}$ tel que
    $V \circ V \subseteq U$ (o\`u $V \circ V = \{(x,z) : \exists y,\, (x,y) \in V,\,
    (y,z) \in V\}$).
\end{enumerate}
(On ne demande \emph{pas} que $U \in \mathcal{U} \Rightarrow U^{-1} \in \mathcal{U}$.)
\end{definition}

\begin{proposition}
Toute quasi-m\'etrique $d$ engendre une quasi-uniformit\'e dont les entourages de base
sont les $U_\varepsilon = \{(x,y) : d(x,y) < \varepsilon\}$.
\end{proposition}

\begin{definition}[Quasi-uniformit\'e de Pervin]
Pour tout espace topologique $(X,\tau)$, la \emph{quasi-uniformit\'e de Pervin} est
engendr\'ee par les entourages~:
\[
  S_U = (U \times U) \cup (U^c \times X), \quad U \in \tau.
\]
Sa topologie associ\'ee est $\tau$.
\end{definition}

\begin{theoreme}[Quasi-uniformisabilit\'e]
Tout espace topologique est quasi-uniformisable (par la quasi-uniformit\'e de Pervin, par
exemple). La cat\'egorie des espaces quasi-uniformes est \'equivalente, en un sens pr\'ecis,
\`a une extension de la topologie g\'en\'erale.
\end{theoreme}

% ------------------------------------------------------------
\section{Compl\'etion de quasi-espaces m\'etriques}
% ------------------------------------------------------------

\begin{definition}[Suite de Cauchy directe]
Dans un espace quasi-m\'etrique $(X,d)$, une suite $(x_n)$ est une \emph{suite de Cauchy
directe} (ou \emph{forward-Cauchy}) si pour tout $\varepsilon > 0$, il existe $N$ tel que
pour tous $N \leq m \leq n$, $d(x_m, x_n) < \varepsilon$.
\end{definition}

\begin{definition}[$y$-limite]
On dit que $(x_n)$ \emph{$d$-converge} vers $\ell$ (ou est \emph{yoneda-convergente}) si
pour tout $\varepsilon > 0$, il existe $N$ tel que pour tout $n \geq N$,
$d(\ell, x_n) < \varepsilon$ \emph{et} pour tout $y \in X$ tel que
$\limsup_{n\to\infty} d(y, x_n) \leq r$, on a $d(y, \ell) \leq r$.
\end{definition}

\begin{definition}[Compl\'etion de Yoneda]
Un espace quasi-m\'etrique est \emph{Yoneda-complet} si toute suite de Cauchy directe
a une $y$-limite. La \emph{compl\'etion de Yoneda} de $(X,d)$ est le plus petit espace
quasi-m\'etrique Yoneda-complet contenant $(X,d)$ comme sous-espace dense.
\end{definition}

\begin{theoreme}[Compl\'etion et domaines]
Soit $(X,d)$ un espace quasi-m\'etrique $T_0$-s\'epar\'e. Si $d$ est \`a valeurs dans
$[0,+\infty)$ et v\'erifie $d(x,y) = 0 \Rightarrow x \leq y$, alors la compl\'etion de
Yoneda de $(X,d)$ est un dcpo muni d'une quasi-m\'etrique compatible avec la topologie de
Scott.
\end{theoreme}

% ------------------------------------------------------------
\section{Continuit\'e et quasi-m\'etriques}
% ------------------------------------------------------------

\begin{definition}[Application non-expansive]
Une application $f\colon (X,d_X) \to (Y,d_Y)$ est \emph{non-expansive} si
$d_Y(f(x),f(y)) \leq d_X(x,y)$ pour tous $x,y$.
\end{definition}

\begin{proposition}
Toute application non-expansive est continue pour les topologies directes.
La r\'eciproque est fausse.
\end{proposition}

\begin{theoreme}[Cat\'egorie des espaces quasi-m\'etriques]
La cat\'egorie $\mathbf{QMet}$ des espaces quasi-m\'etriques $T_0$-s\'epar\'es et
applications non-expansives est compl\`ete et cocompl\`ete. Le foncteur d'oubli
$\mathbf{QMet} \to \mathbf{Top}_0$ pr\'eserve les limites.
\end{theoreme}

% ------------------------------------------------------------
\section{Quasi-m\'etriques et r\'eseaux de Petri (aper\c{c}u)}
% ------------------------------------------------------------

\begin{remarque}
Les quasi-m\'etriques apparaissent naturellement dans l'\'etude des r\'eseaux de Petri et
des syst\`emes de transitions \'etiquet\'es, o\`u la <<~distance~>> de l'\'etat $s$ \`a
l'\'etat $t$ mesure le co\^ut minimal pour passer de $s$ \`a $t$. Cette distance est
g\'en\'eralement non sym\'etrique.
\end{remarque}

% ------------------------------------------------------------
\section{Topologie co-directe et paire conjugu\'ee}
% ------------------------------------------------------------

\begin{definition}[Topologie co-directe]
La \emph{topologie co-directe} (ou topologie conjugu\'ee) est $\tau_{d^{-1}}$, la topologie
engendr\'ee par les boules de la quasi-m\'etrique conjugu\'ee $d^{-1}$.
\end{definition}

\begin{proposition}
Les topologies $\tau_d$ et $\tau_{d^{-1}}$ d\'eterminent un espace bitopologique
$(X, \tau_d, \tau_{d^{-1}})$. La topologie jointe $\tau_d \vee \tau_{d^{-1}}$ co\"incide
avec la topologie de la sym\'etris\'ee $\tau_{d^s}$.
\end{proposition}

\begin{proof}
La topologie jointe est engendr\'ee par les boules ouvertes de $d$ et de $d^{-1}$. Une
boule ouverte de $d^s$ est $B_{d^s}(x,\varepsilon) = B_d(x,\varepsilon) \cap
B_{d^{-1}}(x,\varepsilon)$, ce qui engendre la m\^eme topologie.
\end{proof}

% ------------------------------------------------------------
\section{Exercices}
% ------------------------------------------------------------

\begin{exercice}
Montrer que l'ordre de sp\'ecialisation de la topologie associ\'ee \`a une quasi-m\'etrique
$d$ est donn\'e par $x \leq y \Leftrightarrow d(x,y) = 0$.
\end{exercice}

\begin{exercice}
Soit $(X,d)$ un espace quasi-m\'etrique. Montrer que $d$ est une m\'etrique si et
seulement si $\tau_d$ est $T_1$.
\end{exercice}

\begin{exercice}
Construire une quasi-m\'etrique sur $\{0,1\}$ dont la topologie directe est celle de
Sierpi\'nski.
\end{exercice}

\begin{exercice}
Montrer que la quasi-uniformit\'e de Pervin d'un espace $T_0$ est la plus petite
quasi-uniformit\'e compatible avec la topologie.
\end{exercice}

\begin{exercice}[Compl\'etion de $\Q$]
Soit $d(x,y) = \max(y-x, 0)$ sur $\Q$. D\'ecrire la topologie directe, l'ordre de
sp\'ecialisation, et la compl\'etion de Yoneda de $(\Q, d)$.
\end{exercice}

\begin{exercice}
Montrer que pour un espace quasi-m\'etrique $(X,d)$ \`a valeurs finies, les propri\'et\'es
suivantes sont \'equivalentes~:
\begin{enumerate}
  \item $(X,d)$ est Yoneda-complet.
  \item Tout filtre de Cauchy directe converge.
  \item Toute suite de Cauchy directe a une $y$-limite.
\end{enumerate}
\end{exercice}

\begin{exercice}
Soit $D$ un dcpo avec la topologie de Scott. Construire une quasi-m\'etrique $d$ sur $D$
telle que $\tau_d$ soit la topologie de Scott et $d(x,y) = 0 \Leftrightarrow x \sqsubseteq y$.
(Indication~: utiliser une famille de fonctions Scott-continues \`a valeurs dans
$[0,+\infty)$.)
\end{exercice}
